% Ad Familiares 13.77 --- headnote
% ad-familiares-13-77, late 46 BC (intercalary month), Rome

\textit{Cicero at Rome to Ser. Sulpicius Rufus, proconsul of
Achaia and acclaimed \textit{imperator} by his troops, written
during the first of Caesar's two intercalary months in late 46 BC
(the manuscript dateline: \textit{Scr. Romae mense intere. pr. a.
708}). The opening of the salutation is textually corrupt --- the
Perseus reading is \textit{M. CICERO S. D. n SVLPICIO IMfl}, in
which ``IMfl'' is plainly \textit{IMP.} (imperator) and the stray
``n'' is most likely a misread praenomen abbreviation (Servius =
\textit{Ser.}). The reading here is restored, on the convention
of this cluster's other letters to Sulpicius, to \textit{Cicero
Sulpicio Imp. s.}}

\textit{Three pieces of business in one short letter. First,
Cicero reports that, although he has been keeping away from the
Senate house in these dispirited months, he forced himself to
attend in order to vote in person for the \textit{supplicatio}
(public thanksgiving) decreed on Sulpicius's behalf, and he asks
that Sulpicius's friends be told of his good will so that they
may freely call on him. Second, he recommends M. Bolanus, an old
friend, to Sulpicius's notice. Third --- and the most striking
turn --- he asks Sulpicius to help recover Dionysius, the slave
who had managed Cicero's library, who absconded after pilfering a
considerable number of books and was now passing himself off as
a freedman at Narona on the Illyrian coast. The note ``the matter
in itself is small, but the pain to my mind is great'' is the
most personal sentence in this short note, and one of the few
glimpses in the correspondence of how seriously Cicero took the
violation of his books. The Perseus dateline pins the letter to
an intercalary month --- almost certainly the first of the two
inserted before Caesar's calendar reform took effect; \texttt{meta/works.yaml}
may carry a year-precision placeholder, which is consistent with
the file prefix \texttt{046bc-}.}
